\documentclass{article}
\usepackage{amsmath,amssymb}
\begin{document}
\title{Approximate cancellation in graded shell gluing}
\maketitle

\paragraph{Checked statement.}
Let a finite family of Boolean downward shells be ordered by decreasing target
rank. A shell with center $C$ and deletion depth $d$ consists of all $x\subseteq C$
with $|C|-|x|=d$. For each shell suppose a pointwise limit $F_C$ of $H_C$ reverse
admissible heads is available on the entire deletion cube. Assume $F_C=0$ on
every deletion layer below $d$, and that $F_C$ has the prescribed strict signs
on the entire target shell. Then the union of all shells can be labeled by
exactly $\sum_C H_C$ original-model heads. For a cover of the full Boolean cube,
the same sum bounds the original worst-case quantity.

\paragraph{Why approximate cancellation suffices.}
The space of scores on the finite Boolean cube has its usual finite-dimensional
topology. Write $\mathcal A_H$ for scores of $H$ actual admissible heads. Strict
sign constraints define an open set, so any strictly signed score in
$\overline{\mathcal A_H}$ has the same strict signs as some actual score in
$\mathcal A_H$. Finite weighted sums of independent limits belong to the
closure of the family with the sum of head counts. This follows by continuity
of finite sums and numerator scaling, with no bounds on residues or rates.

For a fixed reverse approximant, center complementation gives strictly
positive denominator slopes. Adding a nonnegative affine penalty, zero on the
center's downset and positive elsewhere, sends the score to zero off that
downset as the penalty grows. Full-cube denominator positivity and slope
admissibility hold for every nonnegative penalty. Taking the original local
limit afterward is legitimate because the closure is closed. This avoids
assuming any uniformity between the two limits.

Finally, any previously protected point in the current downset lies on a
proper deletion layer or on the current target shell. Its new contribution is
therefore zero or correctly signed. Outside the downset the limiting
contribution is zero. Positive finite scaling then gives all required strict
signs, and openness produces an actual original-model score.

\paragraph{Scope.}
The theorem does not assert the existence of an efficient local decoder. In
particular, signs on a sparse subset of a target shell do not suffice: every
target point of that shell must be controlled. This result removes an exact
lower-layer cancellation requirement, not the same-layer interference issue.
The original model, head admissibility, and worst-case definition are unchanged.

\paragraph{Lean declarations.}
All declarations are in
\texttt{HeadComplexity/Results/ApproximateGradedShellTransfer.lean}, namespace
\texttt{HeadComplexity.AxisZeroGluing}. The endpoint is
\texttt{worstCaseHeadComplexity\_le\_approx\_graded\_shell\_cover}.
The accompanying axiom audit lists only
\texttt{propext}, \texttt{Classical.choice}, and \texttt{Quot.sound}.
\end{document}
